\documentclass[11pt,letterpaper]{article}
\usepackage[T1]{fontenc}
\usepackage{lmodern}
\usepackage[margin=1in,headheight=14pt]{geometry}
\usepackage{amsmath,amssymb,amsthm,mathtools}
\usepackage{microtype,booktabs,array,longtable}
\usepackage{enumitem,fancyhdr,xcolor}
\usepackage{listings}
\makeatletter
\let\tableinput\@@input
\makeatother
\usepackage[hidelinks,pdfusetitle]{hyperref}
\usepackage[nameinlink,noabbrev]{cleveref}
\hypersetup{pdfauthor={Research manuscript prepared with ChatGPT},
  pdftitle={Canonical Binomial Decomposition Does Not Preserve Real Roots},
  pdfsubject={Counterexamples to Mu--Welker Question 3.10; exact and asymptotic proofs},
  pdfkeywords={real-rooted polynomials, binomial expansion, recursive decomposition, Newton inequality}}
\setlength{\parskip}{0.3em}
\setlength{\parindent}{1.4em}
\setcounter{tocdepth}{2}
\numberwithin{equation}{section}
\newtheorem{theorem}{Theorem}[section]
\newtheorem{proposition}[theorem]{Proposition}
\newtheorem{lemma}[theorem]{Lemma}
\newtheorem{corollary}[theorem]{Corollary}
\theoremstyle{definition}
\newtheorem{definition}[theorem]{Definition}
\newtheorem{example}[theorem]{Example}
\theoremstyle{remark}
\newtheorem{remark}[theorem]{Remark}
\crefname{theorem}{Theorem}{Theorems}
\crefname{lemma}{Lemma}{Lemmas}
\crefname{proposition}{Proposition}{Propositions}
\crefname{corollary}{Corollary}{Corollaries}
\crefname{definition}{Definition}{Definitions}
\crefname{remark}{Remark}{Remarks}
\crefname{section}{Section}{Sections}
\crefname{table}{Table}{Tables}
\newcommand{\RR}{\mathbb R}
\newcommand{\ZZ}{\mathbb Z}
\newcommand{\NN}{\mathbb N}
\newcommand{\kap}{\kappa}
\newcommand{\Disc}{\operatorname{Disc}}
\DeclareMathOperator{\lk}{lk}
\DeclareMathOperator{\del}{del}
\newcommand{\coeff}[2]{[t^{#1}]#2}
\newcommand{\F}{F}
\newcommand{\G}{G}
\newcommand{\HH}{H}
\pagestyle{fancy}
\fancyhf{}
\fancyhead[L]{\small Canonical binomial decomposition and real roots}
\fancyhead[R]{\small Research manuscript}
\fancyfoot[C]{\thepage}
\renewcommand{\headrulewidth}{0.3pt}
\lstset{basicstyle=\ttfamily\small,columns=fullflexible,keepspaces=true,
  breaklines=true,frame=single,framerule=0.3pt,showstringspaces=false,
  aboveskip=0.8em,belowskip=0.8em,language=Python}
\title{\LARGE\bfseries Canonical Binomial Decomposition\\
Does Not Preserve Real Roots\\[0.5em]
\large Two minimal counterexamples, exact cubic classifications,\\
and explicit failures in every degree}
\author{Research manuscript prepared with ChatGPT}
\date{20 September 2026}
\begin{document}
\maketitle
\thispagestyle{empty}
\begin{abstract}
We give a negative answer to Question~3.10 of Mu and Welker as stated in
arXiv:2503.24076v1. Their canonical recursive decomposition splits a polynomial
with constant coefficient one as $F(t)=G(t)+tH(t)$ using the integer binomial
expansion of each coefficient. Two cubic inputs are exhibited, and each is
minimal in its own sense. The real-rooted polynomial
$F(t)=(1+2t)(1+3t)^2$ has canonical outputs
$G(t)=1+7t+15t^2+8t^3$ and $H(t)=1+6t+10t^2$, with discriminants $-59$
and $-4$; the simple-root polynomial $F(t)=(1+2t)(1+3t)(1+4t)$ has outputs
$G(t)=1+8t+19t^2+11t^3$ and $H(t)=1+7t+13t^2$, with discriminants $-31$ and
$-3$. In both cases \emph{both} outputs have nonreal zeros, and the second
example removes repeated input roots as a possible explanation.

We prove that degree three is minimal. Exhaustive exact finite certificates
show that $8$ is the least linear coefficient among cubic counterexamples and
$9$ the least among cubic counterexamples with distinct input roots. For
$F_m(t)=(1+mt)^3$ we determine all positive integers $m$ for which either
output is real-rooted: $G_m$ precisely for $m\in\{1,2,6\}$, and $H_m$
precisely for $m\in\{1,2,4,6\}$. A separate family gives a complete
classification for the second output when the input has three distinct roots.

The phenomenon is not confined to cubics, and two explicit thresholds are
proved, neither implying the other. For every $d\ge3$ and every integer
$m\ge20(d-1)$, the second output for the repeated-slope input $(1+mt)^d$
violates the first Newton inequality; for every $d\ge3$ and every integer
$n\ge10d^3$, the same happens for the distinct-root consecutive-slope input
$\prod_{j=0}^{d-1}(1+(n+j)t)$. We also derive a two-sided coefficientwise
scaling limit, a general open-region dilation criterion, and the asymmetry
that the first output eventually regains distinct real roots under dilation of
a fixed simple-root input while the second never does. A compressed simplicial
complex realizes the smaller example's full decomposition as a deletion and a
link, while the same input also has a multipartite realization all of whose
links are real-rooted: the failure is local, not a property of the $f$-vector.
Full proofs, integer certificates, and two independent executable
verification stacks are supplied. The Bell--Skandera simplicial-complex
realization question and Mu--Welker Conjecture~3.8 are not resolved by these
counterexamples.
\end{abstract}
\vfill
\noindent\textbf{Status and scope.}
This is an unrefereed research manuscript, not a proof-assistant
formalization. The question and definitions were checked in the
31 March 2025 arXiv version, in the rendered PDF pages rather than only in
search snippets or an experimental HTML conversion. A literature search on
20 September 2026 did not locate an earlier answer, but priority is not
certified and the publisher's final full text was not inspected. The
mathematical claims below refer to the explicitly defined decomposition,
independently of that bibliographic limitation.
\newpage
\tableofcontents
\newpage

\section{The question and the outcome}\label{sec:question}

Let
\begin{equation}\label{eq:F}
 F(t)=1+c_1t+\cdots+c_dt^d,\qquad c_k\in\ZZ_{>0}.
\end{equation}
We call a nonzero real polynomial \emph{real-rooted} if all its complex zeros
are real, allowing repeated zeros. Constant nonzero polynomials are
real-rooted by this convention. A polynomial of the form \eqref{eq:F} cannot
vanish at a nonnegative real number. Thus, if it is real-rooted, it factors
as
\begin{equation}\label{eq:positivefactor}
 F(t)=\prod_{i=1}^d(1+\lambda_i t),\qquad \lambda_i>0.
\end{equation}
The numbers $\lambda_i$ need not be integers. Such a factorization certifies
real-rootedness without any numerical root calculation.

Mu and Welker define a canonical coefficientwise splitting
$F=G+tH$ in \cite[Section~3.2]{MW}. Their Question~3.10 asks whether
real-rootedness of $F$ forces both outputs to be real-rooted. We answer this
question negatively for that precise splitting. The word \emph{recursive},
used for this decomposition in \cite{MW} and in part of the supplementary
material, names this particular arithmetic rule; it does not mean an
arbitrary choice of $G$ and $H$ with $F=G+tH$.

\subsection{Two headline witnesses}

Two cubic inputs carry the negative answer, and both are kept throughout
because they are minimal in two different senses.
The first has the smallest possible linear coefficient:
\begin{equation}\label{eq:maintriple}
\begin{aligned}
 F(t)&=(1+2t)(1+3t)^2=1+8t+21t^2+18t^3,\\
 G(t)&=1+7t+15t^2+8t^3,\qquad \Disc(G)=-59,\\
 H(t)&=1+6t+10t^2,\qquad\quad\ \ \Disc(H)=-4.
\end{aligned}
\end{equation}
The second is the smallest with three \emph{distinct} input roots, and so
removes ``repeated roots caused it'' as an explanation:
\begin{equation}\label{eq:simpletriple}
\begin{aligned}
 F(t)&=(1+2t)(1+3t)(1+4t)=1+9t+26t^2+24t^3,\\
 G(t)&=1+8t+19t^2+11t^3,\qquad \Disc(G)=-31,\\
 H(t)&=1+7t+13t^2,\qquad\quad\ \ \Disc(H)=-3.
\end{aligned}
\end{equation}
Each proof requires only three integer binomial expansions and two negative
discriminants; neither depends on an extensive computation. They are proved
in \cref{thm:main,thm:simple}.

The extensions are substantial. We prove degree-minimality, certify both
minimality statements \emph{within cubics}, classify two infinite cubic
families completely, and give two uniform constructions,
\begin{equation}\label{eq:uniformpreview}
 F(t)=(1+mt)^d\quad(d\ge3,\ m\ge20(d-1)),
 \qquad
 F_{d,n}(t)=\prod_{j=0}^{d-1}\bigl(1+(n+j)t\bigr)\quad(d\ge3,\ n\ge10d^3),
\end{equation}
for both of which $H$ is not real-rooted. We then explain the
large-coefficient mechanism by a limiting transform. In particular, multiple
input roots are not essential, and neither threshold in
\eqref{eq:uniformpreview} implies the other.

\subsection{Results kept in pairs, and what each pair buys}
\label{subsec:pairs}

This manuscript merges two independently written treatments of the same
question (see \cref{subsec:provenance}). Several results come in pairs whose
members look redundant but are not. Each pair is marked at its place in the
text. We list them here with one sentence on what keeping both buys.

\begin{enumerate}[label=P\arabic*.,leftmargin=3.2em]
\item\label{pair:witness}
 \emph{Both cubic witnesses}, \eqref{eq:maintriple} and
 \eqref{eq:simpletriple}, each with its own digit table and discriminant pair
 (\cref{thm:main,thm:simple}): the first is the smallest counterexample,
 and the second shows that a repeated input root is not what causes failure.
\item\label{pair:minimal}
 \emph{Both minimality statements}
 (\cref{prop:minimalcubic,prop:simpleminimal}): $8$ is least among cubic
 counterexamples while $9$ is least among cubic counterexamples with distinct
 roots, so the two numbers answer two different questions and neither
 subsumes the other.
\item\label{pair:threshold}
 \emph{Both all-degree families with their explicit thresholds}
 (\cref{thm:uniform,thm:explicitsimple}): the repeated-slope family has the
 far smaller threshold $20(d-1)$, and the consecutive-slope family is the only
 one for which an explicit threshold is proved with \emph{distinct} roots.
\item\label{pair:limit}
 \emph{Both scaling limits} in \cref{thm:scaling}, on the $H$ side and on the
 $G$ side: only having both exhibits the asymmetry of
 \cref{prop:Grecovery}, which is the sharpest qualitative statement about the
 first output.
\item\label{pair:transfer}
 \emph{Both directions of the relation between the two failures}
 (\cref{lem:transfer} and the counterexample to its converse in
 \eqref{eq:converse}): the lemma alone would misstate the scope, since $G$ can
 be real-rooted while $H$ is not.
\item\label{pair:complex}
 \emph{Both simplicial realizations of the same input} $1+8t+21t^2+18t^3$
 (\cref{subsec:complexes,sec:complex}): the multipartite join has only
 real-rooted links, the compressed complex has a non-real-rooted link, and
 together they prove that the $f$-vector does not determine the
 real-rootedness of every local link.
\item\label{pair:split}
 \emph{Both decompositions of the repeated input}, the canonical one and the
 noncanonical one in \eqref{eq:noncanonical}: the second shows the negative
 answer is about the specified rule and not about the existence of some
 real-rooted splitting.
\item\label{pair:verify}
 \emph{Both verification stacks} (\cref{sec:verification}): one is independent
 inside Python, the other is independent of Python, and two different computer
 algebra systems confirming the same integers is worth more than either alone.
\item\label{pair:cert}
 \emph{Both standalone certificates} (\cref{app:code,app:certificate}): each is
 a hand-checkable disproof on its own, one driven from the search bounds and
 one from the literal digit table.
\item\label{pair:scan}
 \emph{Both cubic enumerations} (\cref{tab:minimal,tab:cubiccounts}): the wider
 scan is the stronger certificate, the narrower one is hand-auditable, and
 their agreement across two different search rectangles is itself a check.
\item\label{pair:tables}
 \emph{Both all-degree witness tables} (\cref{tab:higher} and
 \texttt{results/all\_degree\_certificates.csv}): the dyadic table shows the
 proved thresholds are far from sharp, the certificate file audits a proved
 threshold exactly at and just above its stated value.
\item\label{pair:prov}
 \emph{Both provenance documents} (\texttt{REFERENCES.md} and
 \texttt{STATUS.md}): neither contains all of the page-level source audit, the
 search enumeration, the journal-metadata caveat and the indexing
 correspondence.
\item\label{pair:degtwo}
 \emph{Both rearrangements in the degree-two proof} (\cref{thm:degree} and
 \cref{rem:degtwoalt}): the second exhibits three separately nonnegative
 contributions, which the first hides.
\item\label{pair:k2}
 \emph{Both forms of the $k=2$ estimate} (\cref{lem:k2}): the exact
 triangular-number form drives the transfer lemma, and the two-sided
 inequalities drive every asymptotic estimate; neither substitutes for the
 other.
\item\label{pair:open}
 \emph{Both statements of what is not resolved}
 (\cref{subsec:notresolved}): Bell--Skandera Question~1.1 and Mu--Welker
 Conjecture~3.8 are logically separate, and Proposition~3.9 of \cite{MW} is
 the bridge between them.
\end{enumerate}

\subsection{Provenance}\label{subsec:provenance}

This article is the merger of two research packages prepared on the same day
for the same question, \texttt{mu-welker-binomial-decomposition} and
\texttt{recursive-binomial-real-rootedness}. They proved the shared core
theorem by the same route and differed only in which witness they headlined.
The first package headlined the repeated-root input \eqref{eq:maintriple} and
carried \eqref{eq:simpletriple} as a proposition; the second headlined
\eqref{eq:simpletriple} and carried \eqref{eq:maintriple} in a subsection on
smaller linear coefficients. Both computed the same three integer binomial
expansions, the same two discriminant pairs, the same hockey-stick uniqueness
argument, the same $k=2$ and $k=3$ estimates, the same first Newton
inequality with the same Cauchy--Schwarz proof, the same leading-digit
asymptotic, the same search rectangle for cubics, and the same conclusion that
the Bell--Skandera question and Conjecture~3.8 are untouched. The present text
proves that shared core once. Where the two genuinely diverged --- the
explicit all-degree thresholds --- both results survive in full, in
\cref{thm:uniform,thm:explicitsimple}. The structure, notation
($\kap_k$, $R_k$, $F$, $G$, $H$) and macros of the first package are used
throughout; the second package's $\mathcal S_k$, $\mathcal G_k$, $f$, $g$, $h$
have been rewritten into them. The supplementary files of both packages are
retained side by side, with the second package's verifier renamed to
\texttt{code/verify\_wide.py}.

\subsection{What is not resolved}\label{subsec:notresolved}

The broader problem discussed in \cite[Question~1.1]{MW} asks whether every
real-rooted polynomial \eqref{eq:F} is the face-counting polynomial of a
simplicial complex. That problem, attributed there to Bell and Skandera, is
logically different from Question~3.10, and we do not resolve it: our input
examples already have explicit simplicial-complex realizations, so they cannot
be obstructions to it. Nor do we disprove the coefficientwise inequality
conjectured in \cite[Conjecture~3.8]{MW}; it holds in every example exhibited
below, as recorded explicitly in \cref{subsec:complexes,sec:complex}.
Proposition~3.9 of \cite{MW} is the bridge between those two statements,
relating the coefficientwise conjecture to the realization question; nothing
here bears on that bridge either. These are two independent non-results and
both are stated deliberately.

\section{The canonical binomial splitting}\label{sec:definition}

\subsection{Integer binomial expansions}

For integers $a\ge0$ and $k\ge0$, we use
$\binom ak=0$ when $a<k$, and $\binom00=1$. Negative upper indices are never
needed. For each positive integer $q$ and each positive
integer $k$, there is a unique representation
\begin{equation}\label{eq:binexp}
 q=\binom{a_k}{k}+\binom{a_{k-1}}{k-1}+\cdots+\binom{a_j}{j},
 \qquad a_k>a_{k-1}>\cdots>a_j\ge j\ge1.
\end{equation}
The lower indices are consecutive; the expansion stops when the remainder
becomes zero. We assign zero the empty expansion. Each polynomial coefficient
receives its own expansion: the upper indices in the expansion of $c_k$ are
not shared with those in the expansion of $c_{k+1}$.

For completeness, the greedy construction is elementary. Choose the largest
$a_k$ such that $\binom{a_k}{k}\le q$. If $r$ is the remainder, then
\[
 0\le r<\binom{a_k+1}{k}-\binom{a_k}{k}
       =\binom{a_k}{k-1}.
\]
If $r>0$, the next upper index is therefore less than $a_k$. Continue with
lower index $k-1$. At lower index one, any positive remainder is represented
by one binomial coefficient, so the procedure terminates.

To see uniqueness, fix the first upper index $a$. Every permissible sum of
lower terms is at most
\[
 \sum_{\ell=1}^{k-1}\binom{a-k+\ell}{\ell}
 =\binom{a}{k-1}-1,
\]
by the elementary hockey-stick identity.
Consequently any representation with first term $\binom ak$ lies in
$[\binom ak,\binom{a+1}{k})$. The first term is forced by $q$, and induction
forces every remaining term. This also proves that the integer algorithm
used in the supplementary code implements the specified expansion, and it is
why an exact binary search for the leading index is sufficient.

\subsection{Conventions}\label{subsec:conventions}

All polynomial coefficient lists in the supplementary JSON and CSV files are
in \emph{ascending} power order. A decomposition list may retain zero trailing
coefficients; those are removed before classifying a polynomial's degree, so a
trailing zero is never mistaken for an extra zero of the polynomial. The
source \cite{MW} writes $f_{i-1}$ for the coefficient of $t^i$; we write
$c_i$ for that coefficient and keep the corresponding indexing in both
outputs.

\subsection{The splitting operator}

\begin{definition}\label[definition]{def:split}
For the expansion \eqref{eq:binexp}, define
\begin{equation}\label{eq:KR}
 \kap_k(q)=\sum_{\ell=j}^k\binom{a_\ell-1}{\ell-1},\qquad
 R_k(q)=\sum_{\ell=j}^k\binom{a_\ell-1}{\ell}.
\end{equation}
Set $\kap_k(0)=R_k(0)=0$. For \eqref{eq:F}, set
\begin{equation}\label{eq:GH}
 G_F(t)=1+\sum_{k=1}^dR_k(c_k)t^k,\qquad
 H_F(t)=\sum_{k=1}^d\kap_k(c_k)t^{k-1}.
\end{equation}
\end{definition}

Pascal's identity gives
\begin{equation}\label{eq:pascal}
 R_k(q)+\kap_k(q)=q,\qquad F(t)=G_F(t)+tH_F(t).
\end{equation}
For $q>0$, $\kap_1(q)=1$. Hence both output polynomials have constant
coefficient one. All their coefficients are nonnegative, and $H_F$ has
positive coefficients and degree exactly $d-1$. The polynomial $G_F$ may have
degree smaller than $d$.

These are the outputs denoted $g$ and $h$ in \cite[Section~3.2]{MW}. In
particular, $H_F$ is \emph{not} the usual $h$-polynomial of a simplicial
complex.

\begin{remark}[A warning about other binomial operators]\label[remark]{rem:operator}
The upper index in both sums in \eqref{eq:KR} is lowered by one. Replacing
$\binom{a_\ell-1}{\ell-1}$ by $\binom{a_\ell}{\ell-1}$ would give a different
operator and would not be the question under consideration. The operation is
also not the one that decrements only the \emph{lower} indices, and it is not
obtained by differentiating the polynomial. Concretely, for the simple-root
example \eqref{eq:simpletriple},
\[
 \kap_3(24)=13,\qquad\text{whereas}\qquad
 \tfrac13\coeff{2}{F'(t)}=24 .
\]
Confusing these constructions changes the question. The digit tables in
\cref{sec:counterexample} make the choice of operation explicit, summand by
summand.
\end{remark}

\subsection{Two elementary estimates}

We isolate bounds that will be used repeatedly. \Cref{lem:k2} is kept in both
of its forms; this is pair~\ref{pair:k2}. The exact form drives the
root-location argument of \cref{lem:transfer}, where an inequality would not
suffice; the two-sided inequalities drive every asymptotic estimate in
\cref{sec:everydegree,sec:scaling}, where the exact form is unusable.

\begin{lemma}\label[lemma]{lem:k2}
For every integer $q\ge1$,
\begin{equation}\label{eq:k2exact}
 \kap_2(q)=\left\lceil\frac{\sqrt{8q+1}-1}{2}\right\rceil.
\end{equation}
Equivalently, $\kap_2(q)$ is the smallest positive integer $u$ such that
$q\le u(u+1)/2$; writing $q=\binom n2+r$ with $0\le r<n$, it equals
$n-1+\mathbf 1_{r>0}$. In particular,
\begin{equation}\label{eq:k2bounds}
 \sqrt q\le\kap_2(q)<\sqrt{2q}+1.
\end{equation}
\end{lemma}
\begin{proof}
Write $q=\binom n2+r$, where $0\le r<n$. If $r=0$, then
$\kap_2(q)=n-1$. If $r>0$, the remaining term is $\binom r1$, and
$\kap_2(q)=(n-1)+1=n$. These are exactly the successive triangular-number
intervals described in the statement.

If $u=\kap_2(q)$, then $q\le u(u+1)/2\le u^2$, which proves the lower
bound. For the upper bound, put
$x=(\sqrt{8q+1}-1)/2$. Since $x<\sqrt{2q}$ and
$\lceil x\rceil<x+1$, formula \eqref{eq:k2exact} gives the result.
\end{proof}

\begin{lemma}\label[lemma]{lem:k3}
For every integer $q\ge1$, with $x=(6q)^{1/3}$,
\begin{equation}\label{eq:k3bound}
 \kap_3(q)>\frac{(x-1)(x-2)}2
           =\frac{x^2}{2}-\frac{3x}{2}+1.
\end{equation}
\end{lemma}
\begin{proof}
Let $n$ be maximal with $\binom n3\le q$. Then
\[
 q<\binom{n+1}{3}=\frac{n^3-n}{6}<\frac{n^3}{6},
\]
so $n>x$. The first term in the expansion contributes
$\binom{n-1}{2}$ to $\kap_3(q)$; all other terms contribute nonnegative
integers. The function $(y-1)(y-2)/2$ is strictly increasing for $y>3/2$,
and $x\ge6^{1/3}>3/2$. Hence
\[
 \kap_3(q)\ge\binom{n-1}{2}>\frac{(x-1)(x-2)}2.
 \qedhere
\]
\end{proof}

\section{An exact negative answer}\label{sec:counterexample}

The two theorems of this section are pair~\ref{pair:witness}: the first gives
the smallest counterexample there is, and the second removes the possibility
that a repeated input root is what causes the failure. Neither is a corollary
of the other, and neither is demoted to a remark.

\subsection{The minimal witness}

\begin{theorem}\label{thm:main}
The polynomial $F(t)=(1+2t)(1+3t)^2$ satisfies every input hypothesis of
Mu--Welker Question~3.10 as stated in arXiv:2503.24076v1, but neither output
in its canonical decomposition is real-rooted.
\end{theorem}
\begin{proof}
The input has positive integer coefficients, constant coefficient one, and
zeros $-1/2$ and $-1/3$, the latter of multiplicity two. Its relevant
canonical expansions are
\begin{equation}\label{eq:mainexp}
 8=\binom81,\qquad
 21=\binom72,\qquad
 18=\binom53+\binom42+\binom21.
\end{equation}
The upper indices $5>4>2$ in the last expansion satisfy the strict ordering
requirement. Applying \eqref{eq:KR} gives the following exact table.
\begin{center}
\begin{tabular}{@{}cccc@{}}
\toprule
$k$ & $c_k$ & $R_k(c_k)$ & $\kap_k(c_k)$\\
\midrule
1 & 8 & $7$ & $1$\\
2 & 21 & $\binom62=15$ & $\binom61=6$\\
3 & 18 & $\binom43+\binom32+\binom11=8$
       & $\binom42+\binom31+\binom10=10$\\
\bottomrule
\end{tabular}
\end{center}
Thus the outputs are precisely those in \eqref{eq:maintriple}. The quadratic
$H$ has discriminant $6^2-4\cdot10=-4$, and its zeros are
\begin{equation}\label{eq:Hexactroots}
 \frac{-3+i}{10},\qquad\frac{-3-i}{10}.
\end{equation}
For a cubic $1+at+bt^2+ct^3$, write
\begin{equation}\label{eq:cubicdisc}
 D(a,b,c)=a^2b^2-4b^3-4a^3c-27c^2+18abc.
\end{equation}
This is its discriminant. Direct substitution gives
\[
 D(7,15,8)=11025-13500-10976-1728+15120=-59.
\]
A real cubic with negative discriminant has one real zero and one nonreal
conjugate pair. Hence $G$ also fails to be real-rooted.
\end{proof}

\begin{remark}[Why the discriminant test is exact]\label[remark]{rem:disc}
For a real cubic with leading coefficient $c\ne0$ and roots $r_1,r_2,r_3$,
its discriminant is $c^4\prod_{i<j}(r_i-r_j)^2$. Three distinct real roots
give positive discriminant. A real root together with $u\pm iv$, $v\ne0$,
gives a negative factor $(2iv)^2$ times a positive real factor. A zero
discriminant means repeated roots; a repeated nonreal root and its conjugate
would require degree at least four. Thus a real cubic is real-rooted exactly
when its discriminant is nonnegative. No floating-point approximation is
involved.
\end{remark}

\subsection{The minimal witness with distinct input roots}\label{subsec:simple}

\begin{theorem}\label{thm:simple}
The simple-root polynomial
\[
 F(t)=(1+2t)(1+3t)(1+4t)=1+9t+26t^2+24t^3
\]
has three distinct negative zeros, and neither canonical output is
real-rooted.
\end{theorem}
\begin{proof}
The input zeros are $-1/2$, $-1/3$, $-1/4$, which are distinct. The complete
digit table is \cref{tab:digits}: it displays the two contributions summand by
summand, so the action of \cref{def:split} can be followed term by term. The
coefficient identity is immediate,
\[
 (1+8t+19t^2+11t^3)+t(1+7t+13t^2)=1+9t+26t^2+24t^3 .
\]
The quadratic output has discriminant $7^2-4\cdot13=-3$ and zeros
\begin{equation}\label{eq:simplezeros}
 \frac{-7+i\sqrt3}{26},\qquad\frac{-7-i\sqrt3}{26},
\end{equation}
which alone disproves the assertion that both outputs must be real-rooted.
For the cubic output, \eqref{eq:cubicdisc} gives
\[
 D(8,19,11)=19^2\cdot8^2-4\cdot11\cdot8^3-4\cdot19^3-27\cdot11^2
  +18\cdot11\cdot19\cdot8=-31 .
\]
By \cref{rem:disc} it has one real zero and a nonreal conjugate pair.
\end{proof}

\begin{table}[htbp]
\centering
\renewcommand{\arraystretch}{1.4}
\begin{tabular}{@{}clrr@{}}
\toprule
$k$ & Canonical expansion of $c_k$ & $R_k(c_k)$ & $\kap_k(c_k)$\\
\midrule
1 & $9=\binom91$ & $8$ & $1$\\
2 & $26=\binom72+\binom51$ & $15+4=19$ & $6+1=7$\\
3 & $24=\binom63+\binom32+\binom11$ & $10+1+0=11$ & $10+2+1=13$\\
\bottomrule
\end{tabular}
\caption{The complete arithmetic certificate for the simple-root witness
\eqref{eq:simpletriple}. Every entry is small enough to check by hand.}
\label{tab:digits}
\end{table}

\subsection{What the examples do not contradict}\label{subsec:complexes}

A simplicial complex on a finite vertex set is a family of subsets closed
under taking subsets; its face polynomial is
$\sum_{\sigma}t^{|\sigma|}$, including the empty face. Partition a vertex
set into blocks of sizes $a_1,\ldots,a_d$, and allow faces to contain at most
one vertex from each block. Independent choices in the blocks give the
face polynomial
\begin{equation}\label{eq:join}
 \prod_{i=1}^d(1+a_it).
\end{equation}
Thus the main example, with block sizes $2,3,3$, already has a realization,
as does \eqref{eq:simpletriple} with block sizes $2,3,4$. This construction
also applies to all integer-slope products used below, so every input in this
manuscript is an $f$-polynomial and none of them can obstruct
\cite[Question~1.1]{MW}. This realization is one half of
pair~\ref{pair:complex}; the other half, a differently shaped complex with the
same face counts, is constructed in \cref{sec:complex}.

Even the two outputs in \eqref{eq:maintriple} are face polynomials. For $H$,
take the complete graph on five vertices together with a sixth isolated
vertex, without two-dimensional faces. For $G$, take all fifteen edges on
six vertices, add a seventh isolated vertex, and select any eight of the
ten triangular faces on the first five vertices. These are valid complexes
with the prescribed face counts. The first complex is a subcomplex of the
second. These two realizations serve a different purpose from those of
\cref{sec:complex}: they show that the \emph{outputs} are themselves
$f$-polynomials, so no realizability obstruction is being exhibited on either
side of the splitting.

Their coefficient sequences also satisfy
\[
 (1,6,10,0)\le(1,7,15,8),\qquad
 (1,7,13,0)\le(1,8,19,11)
\]
coordinatewise, after padding the shorter vector with a trailing zero. So
neither example violates the separate inequality in Mu--Welker
Conjecture~3.8. Failure of real-rootedness of $G$ or $H$ is logically
different from failure of these inequalities.

There is also a different, noncanonical split with real-rooted outputs:
\begin{equation}\label{eq:noncanonical}
 (1+2t)(1+3t)^2
 =\underbrace{(1+t)(1+3t)^2}_{1+7t+15t^2+9t^3}
 +t\underbrace{(1+3t)^2}_{1+6t+9t^2}.
\end{equation}
Comparing this with \eqref{eq:maintriple}, the canonical rule transfers just
one unit from the cubic coefficient of the first output to the quadratic
coefficient of the second. Real-rootedness is lost in both. This is
pair~\ref{pair:split}: without \eqref{eq:noncanonical} a reader could not see
that the negative answer concerns the specified canonical rule and not the
existence of some real-rooted decomposition.

\section{Minimality}\label{sec:minimality}

\subsection{No counterexample in degree at most two}

\begin{theorem}\label{thm:degree}
For every real-rooted polynomial of the form \eqref{eq:F} of degree one or
two, both canonical outputs are real-rooted. Therefore three is the minimum
possible degree of a counterexample.
\end{theorem}
\begin{proof}
For $F=1+at$, the outputs are $G=1+(a-1)t$ and $H=1$.

Now let $F=1+at+bt^2$, with $a,b$ positive integers. Real-rootedness means
$a^2\ge4b$, hence $a\ge2\sqrt b$. Put $u=\kap_2(b)$. The outputs are
\[
 G=1+(a-1)t+(b-u)t^2,\qquad H=1+ut.
\]
The second is linear. The first has nonnegative leading coefficient
$b-u=R_2(b)$. If this coefficient is zero, $G$ has degree at most one.
Otherwise, because $2\sqrt b-1\ge1$, its discriminant satisfies
\begin{align*}
 (a-1)^2-4(b-u)
 &\ge(2\sqrt b-1)^2-4b+4u\\
 &=1+4(u-\sqrt b)\ge1,
\end{align*}
where the last inequality uses \cref{lem:k2}. Thus $G$ is real-rooted.
The example in \cref{thm:main} completes the degree-minimality claim.
\end{proof}

\begin{remark}[The same estimate, factored]\label[remark]{rem:degtwoalt}
This is pair~\ref{pair:degtwo}. The same discriminant can be organized as
\begin{equation}\label{eq:degtwoalt}
 (a-1)^2-4(b-u)
 =(a-2\sqrt b)(a+2\sqrt b-2)+4(u-\sqrt b)+1\ \ge\ 1,
\end{equation}
which exhibits three separately nonnegative contributions: the first factor
pair is nonnegative because $a\ge2\sqrt b$ and $b\ge1$, the middle term is
nonnegative by \cref{lem:k2}, and the last is the constant~$1$. The
rearrangement in the proof above hides this decomposition, so both are
recorded.
\end{remark}

\subsection{The smallest linear coefficients among cubics}
\label{subsec:cubicscan}

The following two statements are finite, exact computer-assisted refinements.
Their search bounds are proved here, and the entire enumeration is supplied as
CSV and JSON certificates and as executable code. They are
pair~\ref{pair:minimal}: the first answers ``how small can a cubic
counterexample be?'', the second answers ``how small can one with distinct
input roots be?'', and the two answers are different integers with different
witnesses.

\begin{proposition}\label[proposition]{prop:minimalcubic}
Among real-rooted cubics $1+at+bt^2+ct^3$ with positive integer coefficients,
there is no counterexample with $a\le7$. With $a=8$, the unique
counterexample is $(a,b,c)=(8,21,18)$, the input of \eqref{eq:maintriple}.
\end{proposition}
\begin{proof}
Write the input as $\prod_{i=1}^3(1+\lambda_it)$ with $\lambda_i>0$.
Then
\[
 a=\lambda_1+\lambda_2+\lambda_3,\quad
 b=\lambda_1\lambda_2+\lambda_1\lambda_3+\lambda_2\lambda_3,\quad
 c=\lambda_1\lambda_2\lambda_3.
\]
Since $c\ge1$, arithmetic--geometric mean gives $a\ge3$ and
$b\ge3c^{2/3}\ge3$. Also $b\le a^2/3$ and $c\le a^3/27$.
Thus the complete finite region for $a\le8$ is
\begin{equation}\label{eq:searchbounds}
 3\le a\le8,\qquad
 3\le b\le\lfloor a^2/3\rfloor,\qquad
 1\le c\le\lfloor a^3/27\rfloor.
\end{equation}
For each triple, compute $D(a,b,c)$ from \eqref{eq:cubicdisc}. By
\cref{rem:disc}, retaining exactly those with $D\ge0$ selects exactly the
real-rooted inputs in this region. Compute their canonical outputs by the
greedy rule and test each output by its actual degree and discriminant.
This gives \cref{tab:minimal}. The sole failed output pair is
\eqref{eq:maintriple}. The complete retained list has 124 rows in
\texttt{data/cubics\_a\_le8.csv}; \texttt{verify.py} regenerates it directly.
\end{proof}

\begin{table}[htbp]
\centering
\begin{tabular}{@{}rrrr@{}}
\toprule
$a$ & Candidate triples & Real-rooted inputs & Counterexamples\\
\midrule
3 & 1 & 1 & 0\\
4 & 6 & 2 & 0\\
5 & 24 & 6 & 0\\
6 & 80 & 16 & 0\\
7 & 168 & 33 & 0\\
8 & 342 & 66 & 1\\
\midrule
Total & 621 & 124 & 1\\
\bottomrule
\end{tabular}
\caption{Complete exact search in \eqref{eq:searchbounds}, using the
arithmetic--geometric lower bound $b\ge3$. A counterexample
means that at least one canonical output is not real-rooted. This is the
hand-auditable half of pair~\ref{pair:scan}.}
\label{tab:minimal}
\end{table}

\begin{proposition}\label[proposition]{prop:simpleminimal}
Among real-rooted cubics $1+at+bt^2+ct^3$ with positive integer coefficients
and \emph{distinct} zeros, there is no counterexample with $a\le8$, and with
$a=9$ the input of \eqref{eq:simpletriple}, namely $(a,b,c)=(9,26,24)$, is a
counterexample. Hence $9$ is the least linear coefficient among cubic
counterexamples with distinct input roots, while $8$ is the least without that
restriction.
\end{proposition}
\begin{proof}
Relax the lower bound on $b$ to $b\ge1$ and extend the scan to $a\le14$, which
is legitimate because \eqref{eq:searchbounds} bounds $b$ and $c$ from above
for every $a$ and the additional triples with $b<3$ are merely never
real-rooted. \Cref{tab:cubiccounts} records the resulting exhaustive
enumeration. By \cref{prop:minimalcubic} there is no counterexample at all for
$a\le7$; the unique one at $a=8$ has $D(8,21,18)=0$, hence a repeated zero by
\cref{rem:disc}. At $a=9$ the triple $(9,26,24)$ has $D=4>0$, so its zeros are
distinct, and \cref{thm:simple} shows both outputs fail.
\end{proof}

\begin{table}[htbp]
\centering
\begin{tabular}{@{}rrrr@{}}
\toprule
$a$ & Candidate $(b,c)$ pairs & Real-rooted inputs & Failures\\
\midrule
1--2 & 0 & 0 & 0\\
3 & 3 & 1 & 0\\
4 & 10 & 2 & 0\\
5 & 32 & 6 & 0\\
6 & 96 & 16 & 0\\
7 & 192 & 33 & 0\\
8 & 378 & 66 & 1\\
9 & 729 & 120 & 2\\
10 & 1221 & 203 & 1\\
11 & 1960 & 329 & 0\\
12 & 3072 & 512 & 6\\
13 & 4536 & 762 & 8\\
14 & 6565 & 1109 & 7\\
\midrule
Total & 18794 & 3159 & 25\\
\bottomrule
\end{tabular}
\caption{The wider exact scan, using the weaker lower bound $b\ge1$. A
``failure'' means that at least one output is not real-rooted; the two outputs
need not both fail. These are not counts over all degrees. This is the
stronger half of pair~\ref{pair:scan}; the complete failure list is in
\texttt{results/cubic\_failures.json}.}
\label{tab:cubiccounts}
\end{table}

\begin{remark}[Why both scans are kept]\label[remark]{rem:twoscans}
This is pair~\ref{pair:scan}. \Cref{tab:minimal} uses the rectangle
$3\le b\le\lfloor a^2/3\rfloor$, \cref{tab:cubiccounts} the rectangle
$1\le b\le\lfloor a^2/3\rfloor$; the candidate counts therefore differ
($342$ versus $378$ at $a=8$, for instance) while the real-rooted counts
agree exactly, row by row, everywhere the two overlap. That agreement across
two differently specified search regions, produced by two programs that share
no code, is itself a check. \Cref{tab:minimal} is small enough to audit by
hand against the $124$-row CSV; \cref{tab:cubiccounts} is the stronger
certificate and is the only place where the vanishing of the failure count at
$a=11$, between two nonzero rows, is visible.
\end{remark}

The restriction to cubics is essential to the wording of both
\cref{prop:minimalcubic,prop:simpleminimal}. We have not proved that eight is
the smallest linear coefficient among counterexamples of \emph{all} degrees,
nor that nine is smallest among simple-root counterexamples of all degrees.
The unconditional degree-minimality result is \cref{thm:degree}.

\section{A complete classification for repeated-root cubics}\label{sec:repeated}

For a positive integer $m$, write
\begin{equation}\label{eq:repeated}
\begin{aligned}
 F_m(t)&=(1+mt)^3=1+3mt+3m^2t^2+m^3t^3,\\
 u_m&=\kap_2(3m^2),\qquad v_m=\kap_3(m^3),\\
 H_m(t)&=1+u_mt+v_mt^2,\\
 G_m(t)&=1+(3m-1)t+(3m^2-u_m)t^2+(m^3-v_m)t^3.
\end{aligned}
\end{equation}

\begin{theorem}\label{thm:repeated}
For positive integers $m$,
\begin{align}
 H_m\text{ is real-rooted}&\iff m\in\{1,2,4,6\},\label{eq:Hclass}\\
 G_m\text{ is real-rooted}&\iff m\in\{1,2,6\}.\label{eq:Gclass}
\end{align}
In particular, both outputs have a nonreal conjugate pair for every
$m\ge7$.
\end{theorem}

\subsection{An analytic tail for the quadratic output}

Put $\rho=6^{1/3}$. From \cref{lem:k2,lem:k3},
\[
 u_m<\sqrt6\,m+1,\qquad
 v_m>\frac{\rho^2m^2}{2}-\frac{3\rho m}{2}+1.
\]
It follows that the discriminant $\delta_m=u_m^2-4v_m$ satisfies
\begin{equation}\label{eq:repeatedtail1}
 \delta_m<(6-2\rho^2)m^2+(6\rho+2\sqrt6)m-3.
\end{equation}
The inequalities
\[
 \rho^2>\frac{33}{10},\qquad \rho<\frac{11}{6},\qquad
 \sqrt6<\frac52
\]
are certified by cubing or squaring positive quantities:
$(33/10)^3<36$, $(11/6)^3>6$, and $(5/2)^2>6$.
Therefore
\begin{equation}\label{eq:repeatedtail2}
 \delta_m<-\frac35m^2+16m-3<0\qquad(m\ge27).
\end{equation}
The final quadratic is negative at $m=27$ and strictly decreasing from that
point on. This is a proof for the entire infinite tail, not an extrapolation
from a table.

For $1\le m\le26$, direct integer binomial expansions give
\cref{tab:repeated}. The entries with nonnegative $\delta_m$ occur exactly
at $m=1,2,4,6$. Together with \eqref{eq:repeatedtail2}, this proves
\eqref{eq:Hclass}.

\begin{table}[p]
\centering
\small
\begin{tabular}{@{}rrrrr@{}}
\toprule
$m$ & $u_m$ & $v_m$ & $\Disc(H_m)$ & $D(3m-1,3m^2-u_m,m^3-v_m)$\\
\midrule
\tableinput data/repeated_table.tex
\bottomrule
\end{tabular}
\caption{Exact finite part of the repeated-cubic classification.
At $m=1$, $G_m=(1+t)^2$ has degree two; the last column is the cubic
\emph{formula} with a zero cubic coefficient, not a claim of degree three.
All other $G_m$ in the table have degree three.}
\label{tab:repeated}
\end{table}

\subsection{Why failure transfers to the cubic output}

A root-location argument is more informative than estimating a second
large discriminant.

\begin{lemma}\label[lemma]{lem:transfer}
For $m\ge2$, if $H_m$ is not real-rooted, then $G_m$ is not real-rooted.
\end{lemma}
\begin{proof}
The leading coefficient $v_m$ of $H_m$ is positive. Thus a negative
quadratic discriminant gives $H_m(t)>0$ for every real $t$. For
$-1/m\le t<0$, the identity
\[
 G_m(t)=(1+mt)^3-tH_m(t)
\]
then implies $G_m(t)>0$. For $t\ge0$, nonnegative coefficients and constant
coefficient one also imply $G_m(t)>0$.

For $m\ge2$, $G_m$ genuinely has degree three. Indeed, $m^3\ge8$ makes the
first upper index $n$ in the third binomial expansion at least four, so
$R_3(m^3)\ge\binom{n-1}{3}>0$.

Suppose, for contradiction, that all three roots of $G_m$ are real. They
must all lie strictly to the left of $-1/m$, so
\[
 G_m(t)=\prod_{i=1}^3(1+\lambda_it),\qquad 0<\lambda_i<m.
\]
Let $\varepsilon_i=m-\lambda_i>0$. Comparing linear coefficients gives
$\sum_i\varepsilon_i=1$. Comparing quadratic coefficients now gives
\[
 3m^2-u_m=3m^2-2m+
     \sum_{i<j}\varepsilon_i\varepsilon_j,
\]
so $u_m=2m-\sum_{i<j}\varepsilon_i\varepsilon_j<2m$.
But
\[
 3m^2>\frac{(2m)(2m+1)}2\qquad(m>1),
\]
and the triangular-number characterization in \cref{lem:k2} implies
$u_m>2m$. This contradiction proves the lemma.
\end{proof}

\begin{proof}[Completion of \cref{thm:repeated}]
The classification of $H_m$ and \cref{lem:transfer} exclude every $m$ except
$1,2,4,6$ for $G_m$. At $m=4$, its discriminant is $-31$, so it also fails:
here $H_4$ succeeds while $G_4$ fails. The other three cases can be verified
without a cubic root formula:
\begin{align*}
 G_1(t)&=(1+t)^2,\\
 G_2(t)&=(1+2t)(1+3t+t^2),\\
 G_6(t)&=(1+7t)(1+10t+23t^2).
\end{align*}
The two quadratic factors have positive discriminants $5$ and $8$.
This proves \eqref{eq:Gclass} and the assertion about $m\ge7$.
\end{proof}

The isolated successful values $m=4$ for $H_m$ and $m=6$ for both outputs
show that the small-parameter behavior is not monotone. An infinite-tail
proof is therefore important even when many finite values have been tested.

\section{A complete quadratic-output classification with simple input roots}
\label{sec:simplefamily}

Consider
\begin{equation}\label{eq:simplefamily}
 \widetilde F_m(t)=(1+2mt)(1+3mt)(1+4mt)
                  =1+9mt+26m^2t^2+24m^3t^3.
\end{equation}
All three input roots are distinct for every positive integer $m$. Put
\[
 \widetilde u_m=\kap_2(26m^2),\qquad
 \widetilde v_m=\kap_3(24m^3).
\]
Its quadratic output is
$\widetilde H_m=1+\widetilde u_mt+\widetilde v_mt^2$.

\begin{theorem}\label{thm:simplefamily}
The polynomial $\widetilde H_m$ is real-rooted exactly when $m=3$ or $m=4$.
Thus it has a nonreal conjugate pair for $m=1,2$ and for every $m\ge5$.
\end{theorem}
\begin{proof}
Let $\rho=144^{1/3}$. By \cref{lem:k2,lem:k3},
\[
 \widetilde u_m<\sqrt{52}\,m+1,\qquad
 \widetilde v_m>\frac{\rho^2m^2}{2}-\frac{3\rho m}{2}+1.
\]
Hence
\[
 \Disc(\widetilde H_m)
 <(52-2\rho^2)m^2+(6\rho+2\sqrt{52})m-3.
\]
The rational bounds
\[
 \rho^2>\frac{137}{5},\qquad
 \rho<\frac{21}{4},\qquad
 \sqrt{52}<\frac{29}{4}
\]
follow from $(137/5)^3<144^2$, $(21/4)^3>144$, and $(29/4)^2>52$.
Thus
\begin{equation}\label{eq:simpletail}
 \Disc(\widetilde H_m)<-\frac{14}{5}m^2+46m-3<0
 \qquad(m\ge17).
\end{equation}
Again the last quadratic is negative at the stated endpoint and decreases
thereafter. The finite computations in \cref{tab:simple} cover all
remaining values. Only $m=3,4$ give nonnegative discriminants.
\end{proof}

\begin{table}[htbp]
\centering
\small
\begin{tabular}{@{}rrrr@{}}
\toprule
$m$ & $\widetilde u_m$ & $\widetilde v_m$ & $\Disc(\widetilde H_m)$\\
\midrule
\tableinput data/simple_table.tex
\bottomrule
\end{tabular}
\caption{The finite part of \cref{thm:simplefamily}. The inequality
\eqref{eq:simpletail} covers every $m\ge17$.}
\label{tab:simple}
\end{table}

\subsection{No converse to the transfer lemma}\label{subsec:converse}

There is no general converse to \cref{lem:transfer} outside its particular
repeated-root family. Already at $m=2$ in \eqref{eq:simplefamily},
\begin{equation}\label{eq:converse}
\begin{aligned}
 \widetilde G_2&=1+17t+90t^2+140t^3, &\quad \Disc(\widetilde G_2)&=20>0,\\
 \widetilde H_2&=1+14t+52t^2, &\quad \Disc(\widetilde H_2)&=-12<0.
\end{aligned}
\end{equation}
So the first output can be real-rooted while the second is not. Together with
\cref{lem:transfer} this is pair~\ref{pair:transfer}: keeping only the lemma
would misstate the scope of the relation between the two failures, since it
holds in exactly one direction and only for the repeated-slope cubics. In
\cref{sec:scaling} we prove that $\widetilde G_m$ is real-rooted for all
sufficiently large $m$, even though $\widetilde H_m$ is not. No exact
classification of $\widetilde G_m$ for all $m$ is asserted here.

\subsection{An open region of dilation failures}\label{subsec:region}

\Cref{thm:simplefamily} is one instance of a general criterion, which
identifies an open region of factor data for which the quadratic output must
eventually fail. Its proof uses two tools stated later, the first Newton
inequality \eqref{eq:newton} of \cref{lem:newton} and the asymptotic
\cref{lem:asymptotic}; neither of those depends on anything in this section,
so there is no circularity.

\begin{proposition}\label[proposition]{prop:region}
Let $\lambda_1,\ldots,\lambda_d$ be positive integers with $d\ge3$, and let
$e_2,e_3$ be their second and third elementary symmetric functions. For the
canonical decomposition of
\[
 F_n(t)=\prod_{i=1}^d(1+n\lambda_it),
\]
the output $H_{F_n}$ is not real-rooted for all sufficiently large integers
$n$, provided
\begin{equation}\label{eq:region}
 9(d-1)^3e_3^2>2(d-2)^3e_2^3 .
\end{equation}
\end{proposition}
\begin{proof}
The input coefficients are $e_2n^2$ and $e_3n^3$, so by \cref{lem:asymptotic}
the first two nonconstant coefficients of $H_{F_n}$ satisfy
\[
 b_1=\sqrt{2e_2}\,n+O(1),\qquad
 b_2=\frac{(6e_3)^{2/3}}2n^2+O(n).
\]
Since $H_{F_n}$ has degree $d-1$, its first Newton difference
(\cref{lem:newton} with $n$ there equal to $d-1$) is
\[
 2(d-1)b_2-(d-2)b_1^2
 =\bigl((d-1)(6e_3)^{2/3}-2(d-2)e_2\bigr)n^2+O(n).
\]
The coefficient of $n^2$ is positive exactly when \eqref{eq:region} holds, as
is seen by cubing the two positive quantities $(d-1)(6e_3)^{2/3}$ and
$2(d-2)e_2$. For sufficiently large $n$ the difference is then positive,
contradicting \eqref{eq:newton}.
\end{proof}

For a cubic $1+At+Bt^2+Ct^3$ with $d=3$, so $e_2=B$ and $e_3=C$, criterion
\eqref{eq:region} reduces to
\begin{equation}\label{eq:cubicregion}
 B^3<36C^2 .
\end{equation}
For the simple-root witness \eqref{eq:simpletriple},
$B^3=26^3=17576$ while $36C^2=36\cdot24^2=20736$, so every integer dilation of
that single input eventually yields a counterexample for the second output;
\cref{thm:simplefamily} makes the exceptional set exact in that case.
The criterion is strict, hence stable under sufficiently small perturbations
of the positive factor parameters. Equality in \eqref{eq:region} is not
decided by this leading-order argument.

\section{Uniform constructions in every degree}\label{sec:everydegree}

We now give explicit infinite families in every possible degree. Only the
first three coefficients of the second output are needed.

\subsection{The first Newton inequality}

\begin{lemma}[First Newton inequality]\label[lemma]{lem:newton}
If
\[
 P(t)=1+b_1t+b_2t^2+\cdots+b_nt^n
\]
has positive coefficients, degree $n\ge2$, and only real zeros, then
\begin{equation}\label{eq:newton}
 (n-1)b_1^2\ge2n b_2.
\end{equation}
\end{lemma}
\begin{proof}
All roots are negative, so $P=\prod_{i=1}^n(1+\lambda_i t)$ with
$\lambda_i>0$. By Cauchy--Schwarz,
\[
 \sum_i\lambda_i^2\ge\frac1n\Bigl(\sum_i\lambda_i\Bigr)^2.
\]
Since $b_1=\sum_i\lambda_i$ and
$2b_2=b_1^2-\sum_i\lambda_i^2$, inequality \eqref{eq:newton} follows.
\end{proof}

This is the first of the usual normalized Newton inequalities; the complete
classical sequence is discussed, for example, in \cite{Vondrak}. The
self-contained proof above is all that is needed here.

\begin{remark}[The implication runs one way]\label[remark]{rem:direction}
Failure of \eqref{eq:newton} is sufficient to prove that a polynomial is not
real-rooted; satisfaction of \eqref{eq:newton} is \emph{not} sufficient to
prove that it is. Every use below is in the first direction. This also limits
how the finite Newton certificates of \cref{subsec:margin} may be read: a
positive margin is a complete proof of failure, whereas a nonpositive margin
proves nothing at all.
\end{remark}

\subsection{The repeated-slope family}\label{subsec:repeatedslope}

\begin{theorem}\label{thm:uniform}
For every integer $d\ge3$ and every integer $m\ge20(d-1)$, the canonical
second output $H$ for
\[
 F(t)=(1+mt)^d
\]
is not real-rooted. In particular, the explicit polynomial
\[
 \bigl(1+20(d-1)t\bigr)^d
\]
is a counterexample in every degree $d\ge3$.
\end{theorem}
\begin{proof}
Let
\[
 u=\kap_2\!\left(\binom d2m^2\right),\qquad
 v=\kap_3\!\left(\binom d3m^3\right).
\]
The polynomial $H$ has degree $d-1$, positive coefficients, and initial
terms $1+ut+vt^2$. Its necessary Newton inequality is
\[
 N:=(d-2)u^2-2(d-1)v\ge0.
\]
We will prove the strict reverse inequality. Write
\[
 s=d-1\ge2,\qquad a=\sqrt{d(d-1)},\qquad
 \rho=\bigl(d(d-1)(d-2)\bigr)^{1/3}=(s^3-s)^{1/3}.
\]
The two elementary bounds give
\[
 u<am+1,\qquad
 v>\frac{\rho^2m^2}{2}-\frac{3\rho m}{2}+1.
\]
Therefore
\begin{equation}\label{eq:Nd}
 N<-A m^2+B m-d,
\end{equation}
where
\[
 A=s\bigl(\rho^2-d(d-2)\bigr),\qquad
 B=2(d-2)a+3s\rho.
\]

We establish the elementary estimates $A>s/4$ and $B<5s^2$.
For the first, compare positive cubes:
\begin{align*}
 (s^3-s)^2-\left(s^2-\frac34\right)^3
 &=\frac{16s^4-44s^2+27}{64}\\
 &=\frac{16(s^2-4)^2+84(s^2-4)+107}{64}>0.
\end{align*}
Thus $\rho^2>s^2-3/4$. Since $d(d-2)=s^2-1$, this gives
$\rho^2-d(d-2)>1/4$, as required. For the second estimate, use $a<d$ and
$\rho<s$ to obtain
\[
 B<2d(d-2)+3s^2=5s^2-2<5s^2.
\]
Substituting in \eqref{eq:Nd} yields
\begin{equation}\label{eq:uniformbound}
 N<-\frac{s}{4}m^2+5s^2m-d
   =sm\left(5s-\frac m4\right)-d.
\end{equation}
For $m\ge20s$, the right-hand side is at most $-d<0$. This contradicts
\cref{lem:newton}, completing the proof.
\end{proof}

The constant $20$ is not optimized. It is useful because it gives a short
uniform certificate, without numerical evaluation of radicals or a
parameter-dependent unproved threshold. The exact cubic result of
\cref{thm:repeated} is much stronger in degree three.

\subsection{The consecutive-slope family with distinct roots}
\label{subsec:consecutive}

The inputs of \cref{thm:uniform} have a single repeated root. The next
theorem supplies an explicit threshold for a family with $d$ \emph{distinct}
roots. Neither theorem implies the other: \cref{thm:uniform} has the far
smaller threshold, and \cref{thm:explicitsimple} is the only explicit
threshold proved here for simple-root inputs. This is
pair~\ref{pair:threshold}, and both proofs are given in full.

\begin{theorem}\label{thm:explicitsimple}
For integers $d\ge3$ and $n\ge10d^3$, set
\[
 F_{d,n}(t)=\prod_{j=0}^{d-1}\bigl(1+(n+j)t\bigr).
\]
Then $F_{d,n}$ has positive integer coefficients and $d$ distinct negative
zeros, but its canonical second output $H_{d,n}$ is not real-rooted.
\end{theorem}
\begin{proof}
The input assertions follow from the displayed factorization: each factor
parameter $n+j$ is a positive integer, and the $d$ zeros $-1/(n+j)$ are
distinct. Put
\[
 Q=\sqrt{d(d-1)},\qquad
 R=\bigl(d(d-1)(d-2)\bigr)^{1/3},\qquad
 K=\frac{d-1}{d-2},
\]
and write $H_{d,n}(t)=1+b_1t+b_2t^2+\cdots+b_{d-1}t^{d-1}$, a polynomial of
degree exactly $d-1$ with positive coefficients.

Every product of two factor parameters is at most $(n+d-1)^2$, and every
product of three is at least $n^3$, so the input coefficients $c_2,c_3$ satisfy
\[
 c_2\le\binom d2(n+d-1)^2,\qquad
 c_3\ge\binom d3n^3 .
\]
By \cref{lem:k2}, $b_1=\kap_2(c_2)<\sqrt{2c_2}+1\le Q(n+d-1)+1$. By
\cref{lem:k3} with $x=(6c_3)^{1/3}\ge Rn$, and by monotonicity of
$(y-1)(y-2)$ for $y>3/2$,
\begin{equation}\label{eq:bbounds}
 b_1<Q(n+d-1)+1,\qquad 2b_2>(Rn-1)(Rn-2).
\end{equation}
It therefore suffices to prove
\begin{equation}\label{eq:suffdifference}
 K(Rn-1)(Rn-2)-\bigl(Q(n+d-1)+1\bigr)^2>0,
\end{equation}
since that implies $2Kb_2>b_1^2$, which is the strict reverse of
\eqref{eq:newton} for a polynomial of degree $d-1$.

Expanding the left-hand side of \eqref{eq:suffdifference} gives
\begin{equation}\label{eq:quadraticbound}
 \Delta n^2-B_dn+C_d,\qquad
\begin{aligned}
 \Delta&=KR^2-Q^2,\\
 B_d&=3KR+2Q^2(d-1)+2Q,\\
 C_d&=2K-\bigl(Q(d-1)+1\bigr)^2 .
\end{aligned}
\end{equation}
We use the three uniform bounds
\begin{equation}\label{eq:threebounds}
 \Delta>\frac13,\qquad B_d<3d^3,\qquad C_d>-d^4 .
\end{equation}

For the first, write $u=Q^2=d(d-1)$ and $v=KR^2$. Direct calculation gives
\[
 v^3=\frac{d^2(d-1)^5}{d-2},\qquad u^3=d^3(d-1)^3,\qquad
 v^3-u^3=\frac{d^2(d-1)^3}{d-2}=Ku^2 .
\]
Hence
\[
 v^3-\Bigl(u+\frac13\Bigr)^3
 =Ku^2-u^2-\frac u3-\frac1{27}
 =\frac{u^2}{d-2}-\frac u3-\frac1{27}.
\]
Since $u/(d-2)=d(d-1)/(d-2)>1$, the right-hand side exceeds
$u-u/3-1/27=2u/3-1/27$, and $u\ge6$ makes this positive. Both $v$ and
$u+1/3$ are positive, so taking cube roots gives $v>u+1/3$, which is
$\Delta>1/3$.

For the second bound, $Q<d$, $R<d$ and $K\le2$ give
\[
 B_d<6d+2d(d-1)^2+2d<2d^3+8d\le3d^3\qquad(d\ge3).
\]
For the third, $Q(d-1)+1<d(d-1)+1\le d^2$ and $K>0$ give $C_d>-d^4$.
Thus \eqref{eq:quadraticbound} exceeds
\[
 \frac{n^2}{3}-3d^3n-d^4
 =n\left(\frac n3-3d^3\right)-d^4 .
\]
At $n\ge10d^3$ this is at least $\tfrac{10}3d^6-d^4>0$, and the expression is
increasing in $n$ beyond that point. This proves
\eqref{eq:suffdifference}, so \cref{lem:newton} excludes real-rootedness of
$H_{d,n}$.
\end{proof}

The number $10d^3$ is an explicit sufficient threshold, not an optimized one,
and it is not a claim that $n=10d^3$ is the first failure. The limiting ratio
computed in \cref{subsec:everydegreelimit} approaches one as $d$ increases,
while the finite binomial digits introduce rounding effects; a simple bound
dominating all of those effects is useful even when far from sharp.

\subsection{A compact finite certificate, and how far the bounds are from
sharp}\label{subsec:margin}

For specified $d$ and $n$, the finite certificate is just the integer
\begin{equation}\label{eq:margin}
 M_{d,n}=2(d-1)\kap_3(c_3)-(d-2)\kap_2(c_2)^2 ,
\end{equation}
where $c_2,c_3$ are the quadratic and cubic input coefficients. A positive
value proves failure, by \cref{lem:newton} and \cref{rem:direction}. Neither
the expanded input polynomial nor any zero of any output is needed. For
example, $d=3$ and $n=270$ give
\[
 c_2=220322,\quad c_3=19902240,\quad
 b_1=664,\quad b_2=121112,\quad M_{3,270}=43552 .
\]
The archive records $294$ such checks, at $n\in\{10d^3,\,10d^3+1,\,20d^3\}$
for every $3\le d\le100$, in
\texttt{results/all\_degree\_certificates.csv}. That file audits the proved
threshold of \cref{thm:explicitsimple} exactly at and just above its stated
value.

\Cref{tab:higher} measures something different, and both are kept; this is
pair~\ref{pair:tables}. It records the first successful \emph{dyadic test
values} $m=1,2,4,8,\ldots$ for the repeated-slope inputs $(1+mt)^d$, with last
column the integer $N$ of \cref{thm:uniform}. Each negative entry alone is a
complete certificate of failure for the stated pair $(d,m)$. Already $d=3$
fails at $m=8$, far below the proved threshold $20\cdot2=40$, so the table
shows how far from sharp the proved bound is. The certificate file, by
contrast, says nothing about sharpness and everything about the correctness of
the stated threshold.

\begin{table}[htbp]
\centering
\small
\begin{tabular}{@{}rrrrr@{}}
\toprule
$d$ & $m$ & $u$ & $v$ & $(d-2)u^2-2(d-1)v$\\
\midrule
\tableinput data/higher_degree_table.tex
\bottomrule
\end{tabular}
\caption{Exact Newton-inequality certificates in degrees $3$ through $20$.
The input is $(1+mt)^d$, and $m$ is the first dyadic test value that works,
not necessarily the smallest possible one.}
\label{tab:higher}
\end{table}

\section{Scaling limits and the mechanism of failure}\label{sec:scaling}

The preceding proofs give explicit counterexamples. A limiting description
explains why failure persists and why the two outputs can behave
asymmetrically.

\subsection{Asymptotics of the integer splitting operator}

\begin{lemma}\label[lemma]{lem:asymptotic}
Fix an integer $k\ge2$ and a real number $c>0$. Suppose $q_m$ are positive
integers with
\[
 q_m=cm^k+O(m^{k-1})\qquad(m\longrightarrow\infty).
\]
Then
\begin{equation}\label{eq:kappaasymptotic}
 \kap_k(q_m)=
 \frac{(k!c)^{(k-1)/k}}{(k-1)!}\,m^{k-1}+O(m^{k-2}).
\end{equation}
The implicit constant may depend on $k$, $c$, and the constant in the
hypothesis, but not on $m$.
\end{lemma}
\begin{proof}
Let $a=a(m)$ be the largest upper index in the $k$th binomial expansion of
$q_m$. Then
\[
 \binom ak\le q_m<\binom{a+1}{k}.
\]
For fixed $k$, both bounding binomial coefficients equal
$a^k/k!+O(a^{k-1})$ as $a\to\infty$. First this shows $a=\Theta(m)$, and
then it shows
\[
 a^k=k!c\,m^k+O(m^{k-1}).
\]
Writing $\alpha=(k!c)^{1/k}$ and factoring
$a^k-(\alpha m)^k$ gives $a=\alpha m+O(1)$.

The largest term of $\kap_k(q_m)$ is
\[
 \binom{a-1}{k-1}
 =\frac{\alpha^{k-1}}{(k-1)!}m^{k-1}+O(m^{k-2}).
\]
Every other term has the form $\binom{a_\ell-1}{\ell-1}$ with
$\ell\le k-1$ and $a_\ell<a=O(m)$. There are at most $k-1$ such terms,
and each is $O(m^{k-2})$. Their sum can therefore change only the error
term, proving \eqref{eq:kappaasymptotic}.
\end{proof}

The main term depends only on the leading growth of the input coefficient.
The smaller binomial summands, including the jumps at binomial thresholds,
are one order lower after rescaling. Equivalently, in terms of the argument
itself,
\[
 \kap_k(N)=\frac{(k!N)^{(k-1)/k}}{(k-1)!}+O_k\bigl(N^{(k-2)/k}\bigr)
 \qquad(N\to\infty),
\]
which for $k=2$ is consistent with the exact formula \eqref{eq:k2exact}.

\begin{theorem}\label{thm:scaling}
Fix a polynomial $F(t)=1+\sum_{k=1}^dc_kt^k$ with positive integer
coefficients, and let $F_m(t)=F(mt)$. Define
\begin{equation}\label{eq:QF}
 Q_F(z)=\sum_{k=1}^d
 \frac{(k!c_k)^{(k-1)/k}}{(k-1)!}\,z^{k-1}.
\end{equation}
Then, coefficientwise and uniformly on every compact subset of $\mathbb C$,
\begin{align}
 H_{F_m}(z/m)&=Q_F(z)+O(m^{-1}),\label{eq:Hlimit}\\
 G_{F_m}(z/m)&=F(z)-\frac{z}{m}Q_F(z)+O(m^{-2}).\label{eq:Glimit}
\end{align}
\end{theorem}
\begin{proof}
For $k\ge2$, apply \cref{lem:asymptotic} with $q_m=c_km^k$, and divide by
$m^{k-1}$ to obtain the coefficient of $z^{k-1}$ in
$H_{F_m}(z/m)$. The constant coefficient, coming from $k=1$, is exactly
one, also agreeing with \eqref{eq:QF}. There are finitely many
coefficients, so the coefficientwise errors are uniform on any prescribed
compact set. The identity
\[
 G_{F_m}(z/m)=F(z)-\frac zmH_{F_m}(z/m)
\]
then gives \eqref{eq:Glimit}.
\end{proof}

Both halves of \cref{thm:scaling} are used, and this is pair~\ref{pair:limit}:
\eqref{eq:Hlimit} produces the failures, \eqref{eq:Glimit} produces the
recovery in \cref{prop:Grecovery}, and only the two together exhibit the
asymmetry between the outputs.

\begin{remark}[The limit is not a linear operator]\label[remark]{rem:notlinear}
The map $F\mapsto Q_F$ is not a linear transformation of coefficient vectors:
different coefficients are raised to different fractional powers with
different normalizing constants. Hence the preservation properties of
differentiation, or of other linear root-preserving operators, cannot simply
be transferred to it.
\end{remark}

For the repeated cubic, the limiting quadratic is
\begin{equation}\label{eq:Q3}
 Q_{(1+t)^3}(z)=1+\sqrt6\,z+\frac{6^{2/3}}2z^2.
\end{equation}
Its discriminant is $6-2\cdot6^{2/3}<0$. This predicts eventual failure of
$H_m$; \cref{thm:repeated} supplies the exact exceptional set.

\subsection{The first output recovers simple real roots under dilation}

\begin{proposition}\label[proposition]{prop:Grecovery}
Suppose the fixed polynomial $F$ in \cref{thm:scaling} has $d$ distinct
negative real roots. Then $G_{F_m}$ has $d$ distinct negative real roots for
all sufficiently large integers $m$.
\end{proposition}
\begin{proof}
Choose $d$ disjoint closed intervals on the negative real axis, each
containing one root of $F$ in its interior. Since the roots are simple,
choose the endpoints so that $F$ is nonzero and has opposite signs at the
two endpoints of each interval. Uniform convergence
$G_{F_m}(z/m)\to F(z)$ preserves those endpoint signs for all sufficiently
large $m$. The intermediate value theorem then gives one root of the
rescaled polynomial in each interval.

Also its leading coefficient is $c_d+O(m^{-1})>0$ for sufficiently large
$m$, so its degree is $d$. The $d$ disjoint intervals exhaust its roots,
which are therefore all distinct, negative, and real. Rescaling back does
not change their reality or multiplicities.
\end{proof}

Applied to $F(t)=(1+2t)(1+3t)(1+4t)$, this proves the eventual
real-rootedness of $\widetilde G_m$ asserted in \cref{subsec:converse}.
Meanwhile $\widetilde H_m$ is nonreal for every $m\ge5$. Thus small
perturbation of the input polynomial itself is a good model for $G$, but not
for $H$: the latter converges to the different transform $Q_F$. This
asymmetry --- under dilation the first output eventually becomes real-rooted
while the second never does --- is the most interesting qualitative statement
here about the $G$ side, and it is available only because both limits in
\cref{thm:scaling} were proved.

\subsection{Simple-root failures in every degree}
\label{subsec:everydegreelimit}

For $F(t)=(1+t)^d$, write $Q_d=Q_F$. Its first three coefficients are
\begin{equation}\label{eq:Qdcoeff}
 Q_d(z)=1+A_1z+A_2z^2+\cdots,
 \qquad
 A_1=\sqrt{d(d-1)},\quad
 A_2=\frac{[d(d-1)(d-2)]^{2/3}}2.
\end{equation}
Its degree is $d-1$. The first Newton inequality would require
\begin{equation}\label{eq:Qnewton}
 d(d-1)\ge
 \frac{d-1}{d-2}[d(d-1)(d-2)]^{2/3}.
\end{equation}
Both sides are positive, but the ratio of their cubes is
\begin{equation}\label{eq:ratiolimit}
 \frac{d(d-2)}{(d-1)^2}=1-\frac1{(d-1)^2}<1.
\end{equation}
Therefore \eqref{eq:Qnewton} fails strictly for every $d\ge3$: the limit
object $Q_d$ itself violates the first Newton inequality. This gives a
qualitative limiting explanation of \cref{thm:uniform}, separate from its
explicit bound.

The same limit yields inputs without repeated roots, for which
\cref{thm:explicitsimple} has already supplied an explicit threshold. We
record the limiting statement because it is the conceptual reason, and because
it applies to consecutive slopes shifted by any fixed amount.

\begin{corollary}\label[corollary]{cor:simpleall}
Fix any integer $d\ge3$. For all sufficiently large positive integers $m$,
the polynomial
\begin{equation}\label{eq:consecutiveslopes}
 P_{d,m}(t)=\prod_{j=1}^d\bigl(1+(m+j)t\bigr)
\end{equation}
has positive integer coefficients and distinct negative roots, while its
canonical second output is not real-rooted.
\end{corollary}
\begin{proof}
The input properties are immediate from the distinct positive integer
slopes $m+1,\ldots,m+d$. For fixed $d$, its coefficient of $t^k$ is
$\binom dk m^k+O(m^{k-1})$, because it is the sum of $\binom dk$ products,
each with leading term $m^k$. By \cref{lem:asymptotic}, the coefficient of
$z^{k-1}$ in $H_{P_{d,m}}(z/m)$ tends to the corresponding coefficient of
$Q_d$. The strict reverse of \eqref{eq:Qnewton} is an open inequality in the
first two nonconstant coefficients, so it also holds for $H_{P_{d,m}}(z/m)$
once $m$ is sufficiently large. This polynomial has positive coefficients and
degree $d-1$, so \cref{lem:newton} rules out real-rootedness. Positive
rescaling of the variable preserves whether roots are real.
\end{proof}

The phrase ``sufficiently large'' in \cref{cor:simpleall} is established by a
proved coefficient limit and a strict inequality, and no threshold comes out
of that argument. The explicit threshold for this family is
\cref{thm:explicitsimple}, whose proof is by finite estimates rather than by a
limit. \Cref{thm:uniform} has the different explicit threshold $20(d-1)$ for
the repeated-slope family, which \cref{thm:explicitsimple} does not cover.

\section{Deletion, link, and a compressed realization}\label{sec:complex}

This section reads the identity $F=G+tH$ combinatorially. It explains where a
splitting of this shape comes from, and it shows that the failure exhibited in
\cref{sec:counterexample} is a \emph{local} phenomenon which the $f$-vector
cannot see.

\subsection{The deletion--link identity}

For a vertex $v$ of a simplicial complex $\mathcal C$, the deletion
$\del_v\mathcal C$ comprises the faces not containing $v$, and the link
$\lk_v\mathcal C$ comprises the sets $\tau$ disjoint from $v$ for which
$\tau\cup\{v\}$ is a face. Separating the faces according to whether they
contain $v$ proves
\begin{equation}\label{eq:deletelink}
 f_{\mathcal C}(t)=f_{\del_v\mathcal C}(t)+t\,f_{\lk_v\mathcal C}(t).
\end{equation}
This is the structural shape of $F=G+tH$. The canonical rule of
\cref{def:split} is one arithmetic choice of such a shape, and the next
construction realizes that particular choice geometrically for the witness
\eqref{eq:maintriple}.

\subsection{A compressed complex realizing the full decomposition}

Use vertices $1,\ldots,8$. Include the empty face, all eight singletons, every
pair from $\{1,\ldots,7\}$ as an edge, and the triangles
\begin{equation}\label{eq:trianglefamily}
 \binom{[5]}{3}
 \ \cup\ \bigl\{\{6\}\cup A:A\in\tbinom{[4]}2\bigr\}
 \ \cup\ \bigl\{\{1,5,6\},\{2,5,6\}\bigr\},
\end{equation}
and no faces of size four or more. All subsets of each displayed triangle are
present, so this is a simplicial complex. It has $8$ vertices, $21$ edges and
$10+6+2=18$ triangles, so its face polynomial is exactly the input
$1+8t+21t^2+18t^3$ of \eqref{eq:maintriple}.

The link of vertex $1$ has vertices $2,\ldots,7$; its edges are every pair from
$\{2,3,4,5,6\}$, with vertex $7$ isolated. It has no triangular faces, since
the complex has no four-vertex faces. Its face polynomial is therefore exactly
\[
 1+6t+\binom52t^2=1+6t+10t^2=H(t).
\]
The deletion of vertex $1$ has $7$ vertices, $15$ edges and $8$ triangles, so
its face polynomial is exactly $1+7t+15t^2+8t^3=G(t)$. Direct face counting
thus recovers the complete canonical decomposition \eqref{eq:maintriple}, with
no arithmetic beyond counting.

The triangles in \eqref{eq:trianglefamily} are the first $18$ triples in
\emph{colexicographic} order, where two distinct sets of the same size are
compared by the largest element of their symmetric difference; the edges are
the first $21$ pairs and the singletons the first $8$ elements in the same
order. This finite observation explains why a binomially compressed
realization appears here. No general compression theorem is invoked for this
example.

\subsection{The failure is local, not a property of the \texorpdfstring{$f$}{f}-vector}

The same input $1+8t+21t^2+18t^3$ also has the multipartite realization
\eqref{eq:join} with part sizes $2,3,3$. Every vertex link in that complex is
again a join of discrete sets, so every such link has a real-rooted face
polynomial. The compressed complex above has the \emph{same} face-counting
polynomial, but its link at vertex $1$ is $1+6t+10t^2$, which is not
real-rooted. This is the payoff of pair~\ref{pair:complex}: face counts and
their zeros do not determine the real-rootedness of every local link, so the
failure discovered in \cref{thm:main} is local and cannot be read off the
$f$-vector.

These two realizations are also different in purpose from the ones in
\cref{subsec:complexes}: there the point was that the two \emph{outputs} are
themselves $f$-polynomials, so no realizability obstruction is exhibited on
either side. Here the point is that one $f$-vector supports both a complex all
of whose links are real-rooted and a complex with a link that is not. Neither
fact follows from the other.

Finally, the coefficientwise conjecture is untouched by either realization: as
recorded in \cref{subsec:complexes}, both examples satisfy the coordinatewise
inequalities, so Mu--Welker Conjecture~3.8 is not falsified here.

\section{Reproducibility and the role of computation}\label{sec:verification}

The accompanying files distinguish arithmetic checking from the proofs of
infinite statements. Two verification stacks are supplied, and both are kept;
this is pair~\ref{pair:verify}. They are independent in different senses and
were run on different software.

The core program \texttt{verify.py} requires Python
3.10 or later and only its standard library. It uses arbitrary-precision
integers and exact rational arithmetic. There is no numerical root finder,
floating-point discriminant, or probabilistic test. The second program,
\texttt{code/verify\_wide.py}, shares no code with it: the canonical
expansion, the splitting operator, the discriminants and the search loops are
independently reimplemented, and it carries the wider cubic scan, the
colex-prefix recomputation and the all-degree certificates.

The counterexamples are hand-checkable proofs. Computation is used to
check their arithmetic redundantly, to certify the finite cubic searches in
\cref{prop:minimalcubic,prop:simpleminimal}, and to generate the finite parts
of the exact family classifications. Infinite tails and all-degree conclusions
are proved analytically in the text.

\subsection{Stack one: the core verifier}

\begin{center}
\begin{tabular}{@{}p{0.64\textwidth}r@{}}
\toprule
Exact check & Scope or count\\
\midrule
Two independent greedy implementations compared & 21,105 cases\\
All cubic candidates in \eqref{eq:searchbounds} & 621 triples\\
Retained real-rooted cubic inputs & 124 triples\\
Each of the two cubic families & $1\le m\le1000$\\
Printed higher-degree Newton witnesses & $3\le d\le20$\\
Additional checks at $m=20(d-1)$ & $3\le d\le200$\\
Independent exact Sturm checks (SymPy 1.14.0) & 869 polynomials\\
\bottomrule
\end{tabular}
\end{center}

The two independent greedy implementations use binary search and linear
search, respectively; boundary cases immediately below, at, and above binomial
coefficients are included. All checks remain enabled even when Python is
run with its optimization flag. Invalid negative, nonintegral, or
zero-order inputs are rejected. Output real-rootedness tests remove
trailing zero coefficients before deciding the degree. The extended family
data in \texttt{data/repeated\_cubic\_family.csv} and
\texttt{data/simple\_cubic\_family.csv} cover both cubic families through
$m=1000$ and can be extended with \texttt{--family-limit}.

The supplementary \texttt{minimal\_certificate.py} checks both main
examples independently in a short standard-library script, and
\texttt{independent\_sympy\_check.py} cross-checks the cubic certificate using
exact square-free parts and Sturm root counts on all $621$ input polynomials
and the $248$ outputs of the $124$ admissible inputs, recomputing the
canonical expansions by a separate linear-search routine. That script does not
import \texttt{verify.py} and is not a prerequisite for the proofs.

\subsection{Stack two: the wide-scan verifier}

\begin{center}
\begin{tabular}{@{}p{0.64\textwidth}r@{}}
\toprule
Exact check & Scope or count\\
\midrule
Canonical expansion and Pascal identity, $0\le N\le5000$, $1\le k\le8$
 & 40,008 checks\\
Independent colex-prefix counts & 2,509 checks\\
Real-rooted quadratic inputs, $2\le a\le100$ & 84,575 inputs\\
Candidate cubic triples, $1\le a\le14$ & 18,794 triples\\
All-degree Newton certificates, $3\le d\le100$ & 294 certificates\\
Discriminants recomputed by Sylvester determinant & 6 checks\\
\bottomrule
\end{tabular}
\end{center}

The colex-prefix check is a genuinely different implementation route to the
splitting operator: initial colexicographic segments of the $k$-subsets of
$[12]$ are constructed as explicit lists of subsets, and the number of them
containing vertex $1$ is compared with $\kap_k(N)$, for each $1\le k\le6$ and
every initial segment. The discriminant check recomputes the same six numbers
from the Sylvester determinant of a polynomial and its derivative by exact
rational Gaussian elimination, rather than from the closed formula. The
explicit complex of \cref{sec:complex} is machine-enumerated as a face set,
checked for closure under subsets and under deletion of a vertex, and checked
against the claimed $f$-, deletion- and link-vectors; the result is recorded
in \texttt{results/compressed\_complex.json} with $f$-vector $(1,8,21,18)$,
deletion vector $(1,7,15,8)$ and link vector $(1,6,10)$. The quadratic scan
confirms \cref{thm:degree} computationally over its stated range, which the
proof of course does not require.

The recorded run used Python 3.13.5 and took about $1.28$ seconds; that is an
observation about that run, not a portable performance claim. Both an ordinary
run and a \texttt{python -O} run were executed, and the checks raise rather
than assert so that optimization does not disable them.

\subsection{Two computer algebra systems}

The two optional external checks use different systems, and both are
recorded. The script
\texttt{independent\_sympy\_check.py} uses SymPy~1.14.0, as described above.
Separately, \texttt{code/verify.wl} was evaluated in a connected Wolfram
Language kernel reporting version \texttt{15.0.1 for Linux x86 (64-bit)}; it
returned the coefficient lists of all six polynomials of
\eqref{eq:maintriple} and \eqref{eq:simpletriple}, both decomposition
identities as the zero polynomial, the six discriminants $(4,-31,-3,0,-59,-4)$
for $(F,G,H)$ of the simple-root example followed by the repeated-root one,
real-root counts $(3,1,0)$ for the simple-root triple, the symbolic
cube difference
\[
 \frac{(-2+3d)\,(-1+8d-15d^2+9d^3)}{27\,(-2+d)},
\]
and the answer \texttt{False} to the question whether any real $d\ge3$ makes
that difference nonpositive --- which is the rational positivity step
underlying $\Delta>1/3$ in \cref{thm:explicitsimple}. The exact returned text
is in \texttt{results/wolfram\_output.txt}. Two different systems confirming
the same integers is worth more than either alone, and neither is needed to
run the standard-library verification or to follow the proofs.

These are independent computational implementations, not independent human
refereeing and not a proof-assistant formalization. In particular, the
unbounded theorems are not inferred from finite tests.

\subsection{Running everything}

To regenerate the core certificates and report, run
\begin{lstlisting}[language={}]
python verify.py
\end{lstlisting}
from the archive directory; to extend the finite family checks, run
\begin{lstlisting}[language={}]
python verify.py --family-limit 10000
\end{lstlisting}
The larger range is optional and does not strengthen an already proved
infinite tail. For the second stack,
\begin{lstlisting}[language={}]
python code/verify_wide.py --out reproduced_results
python code/certificate.py
python minimal_certificate.py
\end{lstlisting}
Omitting \texttt{--out} writes to \texttt{results/} and replaces the recorded
data there. The files in \texttt{data/} contain the complete narrow cubic
certificate, extended family data, the higher-degree integer witnesses, and
the small generated tables used by this article; the files in
\texttt{results/} contain the wide cubic scan, the all-degree certificates,
the complex certificate and the recorded external output.

To build the PDF after regenerating the tables, run
\begin{lstlisting}[language={}]
latexmk -pdf -interaction=nonstopmode -halt-on-error article.tex
\end{lstlisting}
Two or more ordinary \texttt{pdflatex} runs can be used instead. The archive
contains the already compiled PDF as well as its source. The bibliography is
inline; \texttt{references.bib} is supplied only for reference-manager import.

\section{Conclusions and precise limits}\label{sec:conclusion}

The proposed preservation principle has a negative answer for the
canonical rule in \cref{def:split}. The smallest-degree obstruction is
already visible in a cubic with eight as its linear coefficient, the smallest
with distinct input roots has nine, and in both cases both outputs fail
simultaneously. The obstruction persists in infinite
families, in every degree at least three, and for inputs with distinct
negative roots, with two explicit thresholds proved by different arguments.
In high degree, the first Newton inequality alone detects
failure; locating nonreal roots numerically is unnecessary.

Several distinctions matter. Both minimal linear coefficients are certified
only among cubics. The complete all-$m$ classifications concern $G_m,H_m$
for $(1+mt)^3$ and the quadratic output for
$(1+2mt)(1+3mt)(1+4mt)$, not every possible cubic family. The threshold
$20(d-1)$ applies to the repeated-slope family $(1+mt)^d$ and the threshold
$10d^3$ to the distinct-root consecutive-slope family; neither covers the
other, and neither is claimed to be sharp. The all-degree theorems prove
failure of $H$, not simultaneous failure of both outputs in every degree; only
the explicit cubic examples prove simultaneous failure. Criterion
\eqref{eq:region} is sufficient, not a classification of every input whose
$H$ fails. None of these claims resolves Bell--Skandera's realization question
or the different coefficientwise conjecture of Mu and Welker, and nothing here
bears on Proposition~3.9 of \cite{MW}, which is the bridge between them.

The limiting transform $F\mapsto Q_F$ suggests further concrete problems:
determine when this transform is real-rooted, sharpen the uniform thresholds
in \cref{thm:uniform,thm:explicitsimple}, determine the exact last exceptional
value of $m$ in each fixed degree, and decide the equality case of
\eqref{eq:region}. These are directions suggested by the present
results, not assertions that an exhaustive literature review has verified
their novelty or open status.

Finally, mathematical validity and priority are separate issues. The
finite counterexamples and the stated proofs can be checked from this
manuscript alone. The literature search did not identify a prior answer,
but this draft makes no unconditional claim of first discovery. Its
bibliographic target is explicitly the versioned question in \cite{MW}.

\appendix
\section{A compact exact search certificate}\label{app:code}

The following pseudocode in executable Python syntax shows the essential
finite enumeration behind \cref{prop:minimalcubic}. Here \texttt{decompose}
implements \cref{def:split}, and \texttt{all\_real\_small} uses the actual
polynomial degree and the exact discriminant. Their complete implementations,
validation, and tests are supplied in \texttt{verify.py}.
\begin{lstlisting}
def D(a, b, c):
    return (a*a*b*b - 4*b*b*b - 4*a*a*a*c
            - 27*c*c + 18*a*b*c)

failures = []
for a in range(3, 9):
    for b in range(3, a*a // 3 + 1):
        for c in range(1, a*a*a // 27 + 1):
            if D(a, b, c) < 0:
                continue
            g, h = decompose([1, a, b, c])
            if not (all_real_small(g) and all_real_small(h)):
                failures.append((a, b, c))

assert failures == [(8, 21, 18)]
\end{lstlisting}
The mathematical bounds \eqref{eq:searchbounds} are what make this a
complete certificate rather than a search for suggestive examples.
Repeated roots are retained by the non-strict input condition $D\ge0$.
Lower-degree outputs are tested separately. Widening the loop to
\texttt{range(1, 15)} with \texttt{b} from~$1$ gives the scan of
\cref{tab:cubiccounts}; that variant is implemented in
\texttt{code/verify\_wide.py}.

\section{A standalone minimal certificate}\label{app:certificate}

This is the other half of pair~\ref{pair:cert}. The program below contains an
entire finite disproof by itself, driven from the literal digit table of
\cref{tab:digits}. It uses only small binomial identities, explicit
coefficient vectors, and the two discriminant formulas, and it depends on no
search program and no third-party package. It raises rather than asserts, so
it survives \texttt{python -O}.
\begin{lstlisting}
from math import comb

def check(condition):
    if not condition:
        raise RuntimeError("Certificate check failed")

digits = [[(9, 1)], [(7, 2), (5, 1)],
          [(6, 3), (3, 2), (1, 1)]]
f, g, h = [1], [1], []
for expansion in digits:
    f.append(sum(comb(a, k) for a, k in expansion))
    g.append(sum(comb(a-1, k) if a-1 >= k else 0
                 for a, k in expansion))
    h.append(sum(comb(a-1, k-1) for a, k in expansion))
check(f == [1, 9, 26, 24])
check(g == [1, 8, 19, 11] and h == [1, 7, 13])
check(all(f[k] == g[k] + h[k-1] for k in range(1, 4)))
check(7**2 - 4*13 == -3)
check(19**2*8**2 - 4*11*8**3 - 4*19**3
      - 27*11**2 + 18*11*19*8 == -31)
print("Exact counterexample verified")
\end{lstlisting}
The factorization $(1+2t)(1+3t)(1+4t)$ certifies the real zeros of the input;
the source definition and the digit table certify that these are the required
outputs, rather than merely some decomposition. The file is
\texttt{code/certificate.py}; the companion
\texttt{minimal\_certificate.py} does the same for \emph{both} witnesses
starting from their factorizations instead of their digit tables.

\section{Source and version audit}\label{app:sources}

The primary source used for the question is the PDF of
arXiv:2503.24076v1, dated 31 March 2025. The splitting is on printed page~5,
Section~3.2; the target question is on printed page~6 and is numbered 3.10.
The PDF page image was checked because the available HTML rendering did not
preserve all displayed theorem numbers correctly. The separate
coefficientwise conjecture is numbered 3.8 and is also on page~6; the
realization question attributed to Bell and Skandera is numbered 1.1;
Proposition~3.9 relates 3.8 to 1.1. The source writes $f_{i-1}$ for the
coefficient of $t^i$; this manuscript writes $c_i$ and keeps the corresponding
output indexing, as stated in \cref{subsec:conventions}.

A publisher listing was found for the same work in \emph{Advances in
Applied Mathematics}, volume~175 (2026), article~103041, 16 pages, with DOI
\texttt{10.1016/j.aam.2026.103041}. The final publisher full text was not
available for inspection during this work; the DOI open attempt returned an
access error. It is therefore not assumed
that the final article retains the same numbering or that it contains no
change to the question. Searches using the title, arXiv identifier, the arXiv
DOI \texttt{10.48550/arXiv.2503.24076}, author names, ``Question 3.10'', and
``recursive decomposition'' with ``real rooted'' and ``counterexample'' did
not locate an earlier resolution; an author's publication list at the
University of Marburg was also consulted. These searches are evidence of a
reasonable check, not proof of worldwide absence of prior work. In particular,
inaccessible journal content, unpublished communications, differently worded
results, or unindexed corrections could contain an earlier answer.

The full source URLs, search scope, and access limitations are recorded in
\texttt{REFERENCES.md} and \texttt{STATUS.md}; both are retained because
neither contains all of the above. No third-party article PDF is redistributed
in the archive. This manuscript restates the necessary definitions and
presents its counterexamples and proofs directly. No checksum manifest is
distributed with this package.

\begin{thebibliography}{9}

\bibitem{MW}
Lili Mu and Volkmar Welker,
\emph{On a question about real rooted polynomials and $f$-polynomials of
simplicial complexes}, arXiv:2503.24076v1, 31 March 2025.
\href{https://arxiv.org/abs/2503.24076v1}{\texttt{arXiv:2503.24076v1}}.
Question~3.10 and Section~3.2, on pp.~5--6, are the version-specific target of
this manuscript. A publisher listing identifies the journal version as
\emph{Advances in Applied Mathematics} \textbf{175} (2026), 103041,
\href{https://doi.org/10.1016/j.aam.2026.103041}{doi:10.1016/j.aam.2026.103041}.

\bibitem{Vondrak}
Jan Vondr\'ak (instructor), Scott Mutchnik (original scribe),
\emph{Lecture~11: Real-rooted Polynomials},
Math~233: Non-constructive methods in combinatorics,
Stanford University, 2018.
\href{https://theory.stanford.edu/~jvondrak/MATH233A-2018/Math233-lec11.pdf}
{Lecture notes on real-rooted polynomials and Newton inequalities}.
Theorem~11.2 there states Newton's inequalities; the specific inequality used
in this article is proved from scratch in \cref{lem:newton}, so nothing is
used as a black box.

\end{thebibliography}
\end{document}
